\chapter{Introduction}
\label{ch:intro}

\algolang\ is a small, statically checked programming language for building AI systems out of cooperating agents. A program declares agents and connects them in flows; the compiler checks the program; the runtime executes it with policies enforced and every agent call recorded.

\section{AlgoLang at a glance}

\begin{algo}[title=brief.algo]
model fast = "claude-haiku-5-5"

agent Writer(facts: text, tone: text = "academic") -> text {
    role  "Technical writer"
    goal  "Write a 200-word brief from {facts} in a {tone} tone"
    model fast
    policy { budget tokens 4000  retry 1 }
}

agent Reviewer(draft: text) -> score {
    role  "Peer reviewer"
    goal  "Score {draft} from 0 to 1 for accuracy"
    policy { require human when confidence < 0.70 }
}

flow main(facts: text) {
    [formal, simple] = par { Writer(facts), Writer(facts, tone = "plain") }
    if Reviewer(formal) < 0.80 {
        formal = Writer(facts, tone = "rigorous academic")
    }
    emit formal, simple
}
\end{algo}

In twenty lines this program declares two agents, runs two writers at the same time, asks a reviewer to score the formal draft, sends any low-confidence score to a human, redrafts when the score is low, and returns both drafts. The compiler has already checked that every name, argument, type and tool is correct before anything runs.

\section{The core idea: agents, not functions}

\begin{definitionbox}{The AlgoLang rule}
There are no functions, methods or lambdas. The only thing a program can call is an \kw{agent}. An agent must state its \kw{goal}; it may state its \kw{role}, the \kw{tools} it \kw{uses}, its \kw{model} and its \kw{policy}. A \kw{flow} orchestrates agents and can do nothing else.
\end{definitionbox}

This rule is enforced by the compiler. Calling anything that is not a declared agent is error \errcode{E101}, and the message explains why:

\begin{terminal}
$ algo check bad_example.algo
✗ bad_example.algo: 8 error(s)
  line 20: E101 'summarise_text' is not an agent - in AlgoLang every unit
           of work is an agent (declare it with `agent summarise_text(...)
           -> type { goal "..." }`)
\end{terminal}

Why insist on this? Because each requirement of an agent answers a question that matters when AI does the work:

\begin{center}
\begin{tabularx}{\linewidth}{@{}L{3.1cm}Y@{}}
\toprule
\thead{Agent clause} & \thead{Question it answers}\\
\goldrule
\kw{goal} & What exactly was this worker asked to do?\\
\kw{role} & In what capacity did it act?\\
\kw{uses} & Which external capabilities could it touch?\\
\kw{model} & Which model produced the result?\\
\kw{policy} & What limits applied: cost, retries, human review, forbidden inputs?\\
audit record & What happened, when, with what confidence?\\
\bottomrule
\end{tabularx}
\end{center}

\section{Design goals}

\begin{enumerate}
  \item \textbf{Agents only.} Work is always visible and accountable.
  \item \textbf{No ambient capabilities.} Tools must be declared and then granted to specific agents.
  \item \textbf{Governance in the language.} Policies are syntax, checked at compile time where possible and enforced at run time.
  \item \textbf{Uncertainty is visible.} Every agent call returns a value \emph{and} a confidence.
  \item \textbf{Accountable by default.} Every agent call is written to an audit log without extra code.
  \item \textbf{Concurrency is a keyword.} \kw{par} and \kw{par for} express multi-agent parallelism directly.
  \item \textbf{Small and familiar.} Brace syntax, about two dozen keywords, five types---easy for people and for coding assistants to read and write.
  \item \textbf{Runs everywhere.} Windows, Linux, macOS, Docker, Jupyter and Google Colab, with an offline backend for teaching and testing.
\end{enumerate}

\section{How AlgoLang compares}

\algolang\ does not replace general-purpose languages. It compiles \emph{to} Python and is meant for the layer where agents are composed and governed.

\begin{center}
\small
\begin{tabularx}{\linewidth}{@{}L{3.6cm}YY@{}}
\toprule
\thead{Concern} & \thead{General-purpose language + agent library} & \thead{AlgoLang}\\
\goldrule
Unit of work & Any function; AI calls hidden inside & Declared agents only (\errcode{E101})\\
Tool access & Whatever the code can import & Declared and granted per agent (\errcode{E104})\\
Human oversight & Written by hand, if at all & \kw{require human when confidence < x}\\
Protected attributes & Convention & \kw{forbid}: compile and run-time checks (\errcode{E108})\\
Cost control & External monitoring & \kw{budget tokens N} per agent\\
Audit trail & Added logging & Automatic, one record per call\\
Parallel agents & Thread or async code & \kw{par} and \kw{par for}\\
Confidence & Rarely exposed & Part of every result\\
\bottomrule
\end{tabularx}
\end{center}

\section{What is in version 0.1}

\begin{center}
\small
\begin{tabularx}{\linewidth}{@{}L{4.2cm}Y@{}}
\toprule
\thead{Component} & \thead{Description}\\
\goldrule
Grammar & Formal grammar in Lark EBNF (Appendix~\ref{app:grammar})\\
Compiler \file{algoc.py} & Parser, semantic checker with 15 error codes, Python code generator\\
Runtime \file{algoruntime.py} & Policy enforcement, concurrency, audit log; mock and live (Claude) backends\\
\code{algo} command & \code{check}, \code{build}, \code{run}, \code{new}, \code{doctor}, \code{notebook}\\
Notebook extension & \code{\%\%algo} cell magic for Jupyter and Google Colab, with diagrams\\
Installers & Windows, Linux/macOS, Docker\\
Editor support & VS Code syntax highlighting and run tasks\\
Examples & Research brief, paper review, error showcase, starter template\\
\bottomrule
\end{tabularx}
\end{center}

\begin{caution}
Version 0.2 Enterprise adds executable tools, modules, organisation policy, 120 domain packs, tamper-evident audit logs, a review queue, a Python SDK and a Prompt UI (Chapter~\ref{ch:enterprise}). Section~\ref{sec:limits} lists every known limitation.
\end{caution}
